% RESEARCH_DRAFT_SHA256: cdca79f26d92be5da4405c52e737d084a38dca08acbaafe0405ffd8bc7b60b67
% Layout and source-reference wrappers expanded for the web renderer.
\documentclass[11pt]{article}
\usepackage[margin=25mm]{geometry}
\usepackage{fontspec,unicode-math}
\setmainfont{Latin Modern Roman}
\setmathfont{Latin Modern Math}
\usepackage{amsmath,amssymb,mathtools,xurl,hyperref}
\hypersetup{colorlinks=true,urlcolor=blue,linkcolor=blue}
\setcounter{secnumdepth}{0}
\setlength{\parindent}{0pt}
\setlength{\parskip}{5pt}
\setlength{\emergencystretch}{3em}
\newcommand{\sref}[2]{\hyperlink{source:#1:#2}{[S#2]}}
\newcommand{\cataloguescope}[1]{\paragraph{Source record.}#1}
\newcommand{\floor}[1]{\left\lfloor#1\right\rfloor}
\begin{document}
\section*{Research-attempt provenance}
This standalone research attempt retains the complete original definition body below,
including its historical source and preparation notes. Those notes describe the
earlier catalogue preparation; the new attempt and its verification are identified
separately after that body. No original ID, catalogue status, or source statement is changed.

\section{Specialized algebra kernels and certified interoperability}\label{statement}

\subsection{Definitions and mathematical statement}

\paragraph{Shared notation.}Standard number systems are $\mathbb N,\mathbb Z,\mathbb Q,\mathbb R,\mathbb C$. Specialized notation and index ranges follow the cited source and the definitions below.


A computer-algebra system implements exact symbolic operations such as polynomial arithmetic, ideal membership, and Gröbner bases. Its computational kernel manages coefficient domains and symbolic data structures. An add-in calls those operations through an interface; interoperability requires matching mathematical objects and conventions across systems.
\paragraph{Record scope and qualifications.}
The catalogue title is retained as metadata. Named certificates, contracts, bounds or frameworks in that title are not asserted as results of the original source extract.
This is a software-architecture discussion with a participant response about kernel access. It contains no complexity theorem or specified implementation to certify. The supplied catalogue marks this record partially solved; its present research status was not independently checked.
The source text below is rendered literally where necessary; damaged formulas and composite extracts are not silently repaired into a new mathematical claim.
\paragraph{Statement recorded in the supplied source.}
Question 69 (Roehrig). Why are there so many different systems for computational alge-bra? Couldn't they be add-ins to other systems? Stillman: Not likely to work since those other systems (Mathematica, Maple, etc.) do not give access to their kernels. \href{https://aimath.org/WWN/compalgstat/compalgstat.pdf}{[S2]}
\paragraph{Source status.}
The supplied catalogue labels this source record partially solved; that metadata was checked against the supplied snapshot on 2026-10-11. The label is preserved for provenance and is not an independent verification of present open status or mathematical correctness.
\subsection{Short English statement}
The source discusses whether specialized algebra systems could be integrated into other software.
\subsection{Sources}
\begin{itemize}
\item[S1] Ulam AI and the UnsolvedMath Contributors, UnsolvedMath: A Curated Collection of Open Mathematics Problems, catalogue entry. Catalogue curation and metadata: CC BY 4.0. The source extract is retained; definitions, English summaries and qualifications are editorial additions.
\url{https://huggingface.co/datasets/ulamai/UnsolvedMath}.
\item[S2] AIM Workshop Problem Lists, source linked by the catalogue.
\url{https://aimath.org/WWN/compalgstat/compalgstat.pdf}.
\end{itemize}
\paragraph{Preparation and review.}
The original source and UnsolvedMath contributors retain attribution for the source extract. Codex (GPT-6; exact serving revision not exposed) assisted with definitions, English summaries, scope notes and formatting. No new human mathematical authorship or independently verified solution is asserted. Catalogue identity checked against the complete supplied registry and current repository. Distinct upstream records, including repeated questions and broad or damaged extracts, are retained for complete coverage of the source CS category. A repeated record is not a new independent conjecture. Complete mathematical-equivalence and current-literature review has not been performed. Preparation and final review follow contribution prompts C and D, with source gaps disclosed. Structural and rendering checks do not establish a new theorem or open status.
\label{end:statement}

\section{Literature and context check}
Question 69 and its response are a historical architecture discussion in the AIM source. They do not establish the capabilities of present software. The selected route tests whether a caller can verify one algebraic service from exact output certificates, without access to the service's implementation.
\section{Research attempt: A verifiable algebra service without trusting its internal kernel}
Consider a service claiming to compute the monic greatest common divisor of two polynomials $f,g\in\mathbb Q[x]$, not both zero. Require output polynomials $d,u,v,a,b$ with $d$ monic and
\[
f=ad,\qquad g=bd,\qquad d=uf+vg.
\]
The caller can check these identities by exact polynomial multiplication and comparison. Acceptance proves the claimed answer: the first two equations show $d$ divides both inputs. If $h$ divides both inputs, write $f=hf_1$, $g=hg_1$; the last equation becomes $d=h(uf_1+vg_1)$, so $h$ divides $d$. These two properties define a greatest common divisor up to a nonzero scalar, and the monic condition fixes that scalar. The certificate proof uses no internal operation of the backend.

Both halves of the check are necessary. A divisibility-only backend can always return $d=1$, which divides $f=x^2$ and $g=x^3$ but is not their monic greatest common divisor because the common divisor $x^2$ does not divide one. Conversely an ideal-membership-only backend can return $d=x^2+x^3=f+g$. This is a valid linear combination but does not divide $f$: if $x^2+x^3$ divided the nonzero degree-two polynomial $x^2$, the product-degree formula would force a polynomial quotient of negative degree. Thus one-sided certificates can accept false results.

For an explicit valid response on that instance, take $d=x^2$, $a=1$, $b=x$, $u=1$, $v=0$. All three identities hold. This shows a precise interface in which trust can be confined to elementary exact arithmetic. It does not show that every difficult algebraic result has a short certificate, nor that certificate generation is cheap. Verification must charge the sizes of all supplied coefficients and polynomials, including possibly large intermediate products.
\section{Verification and adversarial audit}
The proof covers one zero input as long as the pair is not both zero; then the certificate still forces the correct nonzero monic divisor. Both-zero inputs are excluded because the monic normalization is unsuitable for the usual zero gcd. The two false responses are rejected by different required identities.
\paragraph{Audit corrections.} A historical assertion about unavailable kernel access is not repeated as a current fact. Existence of a checkable gcd certificate is not generalized to all algebraic computations.
\paragraph{External computation.} None; symbolic arguments and boundary checks only.
\section{Scope and remaining gaps}
This implements a mathematical contract for one standard service, not a software ecosystem or interface to the historical named systems. Kernel-access claims are not updated without current evidence; implementation, interoperability, and performance requirements in the original discussion remain outside this proof.
\paragraph{Integration review.} Codex (GPT-6; exact serving revision not exposed) reviewed the written arguments and expanded layout and source-reference wrappers. Repository metadata and the production web renderer were checked using Python 3.13 and JavaScript with MathJax 3.2.2. This model-assisted review is not independent human proof verification; no PDF compilation or formal proof checking is asserted.

\section{Completion Estimate}
\noindent\textbf{COMPLETION: 0\%}
The calculations below the preserved source statement investigate a concrete route,
but establish no substantiated new progress toward the entire original target.
Known special cases and failed sufficient conditions are not credited as a solution.
\section{References}
The original source references and catalogue attribution are preserved in the
Sources section above. The following URLs were actually consulted for this attempt;
retention of an additional original reference does not claim fresh consultation of it.
\begin{itemize}
\item \url{https://aimath.org/WWN/compalgstat/compalgstat.pdf} (consulted 2026-10-11). 
\end{itemize}
\section{Model and contributor disclosure}
OpenAI gpt-6-astra, requested reasoning configuration ultra, generated this research
attempt on 2026-10-11. The exact serving revision is not exposed. Codex GPT-6
performed orchestration, repository-format review, and a model-assisted mathematical
audit of the written arguments; its exact serving revision
is also not exposed. The mathematical route, boundary checks and final adversarial
self-audit followed contribution prompts A then B in
\texttt{docs/contribution\_prompts.md}, after reading \texttt{CONTRIBUTING.md}.
Python 3.14.2 assembled the standalone documents; any mathematical computation is
specified in Verification. No theorem prover was used. No human mathematical
authorship, substantial human proof correction, or independent human proof checking
is asserted. Model-assisted review, source retrieval, and structural checks do not
establish correctness or novelty. Full private reasoning transcripts are unavailable
as reproducibility artifacts; the complete arguments, inputs and outcomes actually
used are stated here or in the companion batch audit.
\section{Final status}
Unresolved as a new research contribution. This is a concrete route investigation,
not a claim of new complete resolution or a current-status certification. A known
answer reproduced for context does not become a new contribution. The catalogue's
source-status label is preserved; this attempt status does not assert that a
source-labelled solved question is open.
\end{document}
